% GENERATED FILE. Edit canonical lessons, then run npm run generate:books.
% canonical-source-hash: fnv1a64:f8e103d4e60a818b
% canonical-lessons: FR-C361-stresse, FR-C361-epuise, FR-C361-efficace, FR-C361-efficacite, FR-C361-productivite

\chapter{Trop de travail}
\label{ch:fr-trop-de-travail}
\hlchaptermodality{361}

\begin{chapteropening}
\textbf{By the end of this chapter:} \emph{I can say I am stressed or worn out and talk about being effective, efficiency and productivity.}

Complete the last lesson of chapter 361: I can say I am stressed or worn out and talk about being effective, efficiency and productivity.
\end{chapteropening}

\section[stressé]{stressé — stressed}

\label{lesson:FR-C361-stresse}



\begin{quote}
Before the new one: say the French for the skill, then the French for the training.
\end{quote}

\subsection*{You'll want to know: stressé}
\textbf{stressé} — \textquotedblleft{}stressed\textquotedblright{}.

\begin{quote}
Je suis stressé à cause du travail.
\end{quote}

\begin{quote}
\emph{I am stressed because of work.}
\end{quote}

Feminine \emph{stressée}. Built on English \emph{stress}.

\begin{cousinweb}
\emph{stressé} is built on \emph{stress}, which was borrowed from English \emph{stress} \textquotedblleft{}strain, pressure\textquotedblright{} (EtymDB). The English word goes back to Middle English \emph{destresse} \textquotedblleft{}distress\textquotedblright{}.
\end{cousinweb}

\begin{grammarlens}[title={What you've built}]
One more word for telling, explaining and describing at B1.
\end{grammarlens}

\subsection*{Your turn}
\begin{itemize}
\raggedright
  \item \emph{Say it:} \emph{stressé}
  \item \emph{Say it:} \emph{stressé}, once more
  \item \emph{Say it:} a sentence of your own with \emph{stressé}
  \item \emph{Recall:} say the French for the skill, then the French for the training, then say \emph{stressé} again
\end{itemize}

\begin{tcolorbox}[breakable,colback=teal!4,colframe=teal!35!black,title={Before you move on}]
What does stressé mean? (\textquotedblleft{}Stressed\textquotedblright{}.) Say it once more.
\end{tcolorbox}
\section[épuisé]{épuisé — exhausted, worn out; sold out}

\label{lesson:FR-C361-epuise}



\begin{quote}
Before the new one: say the French for the training, then the French for stressed.
\end{quote}

\subsection*{You'll want to know: épuisé}
\textbf{épuisé} — \textquotedblleft{}exhausted, worn out; sold out\textquotedblright{}.

\begin{quote}
Après ce voyage, je suis épuisé.
\end{quote}

\begin{quote}
\emph{After that trip, I am exhausted.}
\end{quote}

Of a book or a product: sold out.

\begin{cousinweb}
\emph{épuisé} is the participle of \emph{épuiser} \textquotedblleft{}to drain dry\textquotedblright{}, built on \emph{puits} \textquotedblleft{}a well\textquotedblright{}, which comes from Latin \emph{puteus} \textquotedblleft{}a well\textquotedblright{} (Lewis \& Short): drained like a well.
\end{cousinweb}

\begin{grammarlens}[title={What you've built}]
One more word for telling, explaining and describing at B1.
\end{grammarlens}

\subsection*{Your turn}
\begin{itemize}
\raggedright
  \item \emph{Say it:} \emph{épuisé}
  \item \emph{Say it:} \emph{épuisé}, once more
  \item \emph{Say it:} a sentence of your own with \emph{épuisé}
  \item \emph{Recall:} say the French for the training, then the French for stressed, then say \emph{épuisé} again
\end{itemize}

\begin{tcolorbox}[breakable,colback=teal!4,colframe=teal!35!black,title={Before you move on}]
What does épuisé mean? (\textquotedblleft{}Exhausted, worn out; sold out\textquotedblright{}.) Say it once more.
\end{tcolorbox}
\section[efficace]{efficace — effective, efficient}

\label{lesson:FR-C361-efficace}



\begin{quote}
Before the new one: say the French for stressed, then the French for exhausted.
\end{quote}

\subsection*{You'll want to know: efficace}
\textbf{efficace} — \textquotedblleft{}effective, efficient\textquotedblright{}.

\begin{quote}
Ce médicament est très efficace.
\end{quote}

\begin{quote}
\emph{This medicine is very effective.}
\end{quote}

One form for both genders.

\begin{cousinweb}
\emph{efficace} comes from Latin \emph{efficāx} \textquotedblleft{}effective\textquotedblright{} (EtymDB; Lewis \& Short), from \emph{efficere} \textquotedblleft{}to bring about\textquotedblright{}, built on \emph{facere} \textquotedblleft{}to do, to make\textquotedblright{}.
\end{cousinweb}

\begin{grammarlens}[title={What you've built}]
One more word for telling, explaining and describing at B1.
\end{grammarlens}

\subsection*{Your turn}
\begin{itemize}
\raggedright
  \item \emph{Say it:} \emph{efficace}
  \item \emph{Say it:} \emph{efficace}, once more
  \item \emph{Say it:} a sentence of your own with \emph{efficace}
  \item \emph{Recall:} say the French for stressed, then the French for exhausted, then say \emph{efficace} again
\end{itemize}

\begin{tcolorbox}[breakable,colback=teal!4,colframe=teal!35!black,title={Before you move on}]
What does efficace mean? (\textquotedblleft{}Effective, efficient\textquotedblright{}.) Say it once more.
\end{tcolorbox}
\section[l'efficacité]{l'efficacité — the efficiency, effectiveness}

\label{lesson:FR-C361-efficacite}



\begin{quote}
Before the new one: say the French for exhausted, then the French for effective.
\end{quote}

\subsection*{You'll want to know: l'efficacité}
\textbf{l'efficacité} — \textquotedblleft{}the efficiency, effectiveness\textquotedblright{}.

\begin{quote}
Il travaille avec beaucoup d'efficacité.
\end{quote}

\begin{quote}
\emph{He works very efficiently.}
\end{quote}

Feminine.

\begin{cousinweb}
\emph{efficacité} is a learned borrowing from Latin \emph{efficācitās} \textquotedblleft{}efficacy\textquotedblright{} (Lewis \& Short), from \emph{efficāx} \textquotedblleft{}effective\textquotedblright{}, built on \emph{facere} \textquotedblleft{}to do\textquotedblright{}.
\end{cousinweb}

\begin{grammarlens}[title={What you've built}]
One more word for telling, explaining and describing at B1.
\end{grammarlens}

\subsection*{Your turn}
\begin{itemize}
\raggedright
  \item \emph{Say it:} \emph{l'efficacité}
  \item \emph{Say it:} \emph{l'efficacité}, once more
  \item \emph{Say it:} a sentence of your own with \emph{l'efficacité}
  \item \emph{Recall:} say the French for exhausted, then the French for effective, then say \emph{l'efficacité} again
\end{itemize}

\begin{tcolorbox}[breakable,colback=teal!4,colframe=teal!35!black,title={Before you move on}]
What does l'efficacité mean? (\textquotedblleft{}The efficiency, effectiveness\textquotedblright{}.) Say it once more.
\end{tcolorbox}
\section[la productivité]{la productivité — the productivity}

\label{lesson:FR-C361-productivite}



\begin{quote}
Before the new one: say the French for effective, then the French for the efficiency.
\end{quote}

\subsection*{You'll want to know: la productivité}
\textbf{la productivité} — \textquotedblleft{}the productivity\textquotedblright{}.

\begin{quote}
La productivité de l'usine a augmenté.
\end{quote}

\begin{quote}
\emph{The factory's productivity has gone up.}
\end{quote}

From \emph{productif}, productive.

\begin{cousinweb}
\emph{productivité} is built from \emph{productif} and \emph{-ité}. The family goes back to Latin \emph{prōdūcere} \textquotedblleft{}to lead forth, to bring out\textquotedblright{}, built on \emph{dūcere} \textquotedblleft{}to lead\textquotedblright{} (Lewis \& Short).
\end{cousinweb}

\begin{grammarlens}[title={What you've built}]
One more word for telling, explaining and describing at B1.
\end{grammarlens}

\subsection*{Your turn}
\begin{itemize}
\raggedright
  \item \emph{Say it:} \emph{la productivité}
  \item \emph{Say it:} \emph{la productivité}, once more
  \item \emph{Say it:} a sentence of your own with \emph{la productivité}
  \item \emph{Recall:} say the French for effective, then the French for the efficiency, then say \emph{la productivité} again
\end{itemize}

\begin{tcolorbox}[breakable,colback=teal!4,colframe=teal!35!black,title={Before you move on}]
What does la productivité mean? (\textquotedblleft{}The productivity\textquotedblright{}.) Say it once more.
\end{tcolorbox}
